\section{Derivation of Adiabatic Time-Evolution Operator}\label{sec:kato_appendix}

In this section we will justify the form of the adiabatic Hamiltonian,
\begin{align}
    K(t) = i\left[ \partial_t P_I(t), P_I(t) \right],
\end{align}
as the generator of the correct adiabatic time-evolution for a given initial projector, $P_0$. We introduce a parameter $T$ that defines the period over which the system evolves, $t \in [0,T]$. Changing variables to $s = tT^{-1}$, we arrive at the scaled Schr\"odinger equation
\begin{align}\label{eqn:scaled_schrodinger}
    i \partial_s U_T(s) = TH(s) U_T(s),
\end{align}
where we can use $T$ to parametrise how slowly the system changes. We will also define the scaled adiabatic Hamiltonian
\begin{align}
    \bar K(s) = i[\partial_s P_I(s), P_I(s)].
\end{align}
Since $\partial_s = T \partial_t$, so this quantity is related to $K$ by $\bar K = T K$.\par  
Our aim will be to show that in the limit of large $T$, the time-evolved state $P(s) = U_T(s) P_0 U^{\dag}_T(s)$, can be approximated by the instantaneous projector, $P_I(s) = \sum_{i \in \textup{band}} \ket{\psi_i(s)} \bra{\psi_i(s)}$, where $\ket{\psi_i(s)}$ are eigenstates of $H(s)$. We wish to define an adiabatic time-evolution operator $U_A(s)$ such that
\begin{align}\label{eqn:U_A_def}
     P_I(s) = U_A(s) P_0 U_A^\dag(s).
\end{align}
We propose an ansatz for $U_A$ that satisfies the following conditions,
\begin{align}
    U_A(0) &= \1, \\
    i \partial_s U_A(s) &= \left [ TH(s) + \bar K(s) \right ] U_A(s).
\end{align}
To verify that $ U_A$ produces the right dynamics, let us consider the quantity
\begin{align}
    G(s) =  U_A(s)^\dag P_I(s) U_A(s).
\end{align}
In the interest of keeping equations legible, we will drop explicit $s$-dependence from our operators. From hereon in it is safe to assume that, unless stated otherwise, all operators are being evaluated at time $s$. We take the $s$-derivative of $ G$,
\begin{align}
    \partial_s  G &= \partial_s \left [ ( U_A^\dag  P_I)( P_I  U_A) \right ].
\end{align}
Using the identity
\begin{align}
    \partial_s ( P_I  U_A) = \partial_s P_I  U_A -i  P_I \left ( T H(s) + \bar K(s)  \right )  U_A,
\end{align}
we can show that
\begin{align}
    \partial_s  G &= 
    U_A^\dag \left ( iT [H,P_I] + \{P_I,\partial_s P_I\}  + i [\bar K,P_I] \right )U_A.
\end{align}
The first term in this expression vanishes, since $P$ and $H$ commute by definition. Next we can use the identities $\partial_s P_I = \{P_I, \partial_s P_I\}$ and $P_I \partial_s P_I P_I = 0$ to show that the second and third terms respectively equal $\partial_s P_I$ and $-\partial_s P_I$, cancelling out. Thus we have shown that $\partial_s G = 0$ for all $s$. Since $G(0) = P_0$, we see that
\begin{align}
    G(s) = P_0, \  \forall\, s,
\end{align}
confirming that our definition satisfies eqn.~\ref{eqn:U_A_def}.\par
Now all that remains is to show that, as $ T \rightarrow \infty$, the actual dynamics, $U_T$, asymptotically approach the adiabatic dynamics described by $U_A$. We will omit the proof of this result, which is simply a proof of the adiabatic theorem, however a complete justification can be found in~\cite{avron_adiabatic_1987}. We quote the result, which compares the projector generated by non-adiabatic dynamics, $P(s) = U_T P_0 U^{\dag}_T$ compared to that generated by the adiabatic time evolution $P_I(s) = U_A P_0 U^{\dag}_A$. The result states that
\begin{align}
    |P(s) - P_I(s)| = \mathcal O (T^{-1})
\end{align}
Thus we can see that in the limit of large $T$, $U_T$ generates the correct time evolution.
